\chapter{附录}\label{ch:appendix}

\section{特殊内核功能的实现}\label{sec:implementation-of-special-kernel-features}

在RMP中，有一部分功能没有被内核默认实现。但是在RMP中实现这些功能是可能的；具体的实现原理如下所列，如果用户需要这些功能则要自行实现。

\subsection{无节拍内核的实现}\label{subsec:implementation-of-a-tickless-kernel}

\protect\phantomsection\label{subsec:implementation-of-a-tickless-kernel-2}无节拍内核要求系统具备一个大时长、高精度的定时器，并且由该定时器产生可变的调度器时间中断。时间片和定时器填充值之间的比例可以由实现者按需自由决定。为了方便实现无节拍内核，系统提供两个函数\_RMP\_Tim\_Future和\_RMP\_Tim\_Elapse；如果不实现无节拍内核可以不必理会它们。

为了使定时器在最近的潜在超时时点产生中断，在每次发生定时器中断\footnote{如定时器中断钩子RMP\_TIM\_HOOK；定时器中断可能意味着（1）某个线程耗尽了自己的时间片，（2）某个线程延时的时间已到，（3）硬件定时器的最长可配置超时时间不够长而需要将一个定时分多段完成；在情况（1）中一定会发生上下文切换，在情况（2）中可能会发生上下文切换，在情况（3）中不可能发生上下文切换；在检测到上下文切换悬起时，不需要进行此次定时器编程，因为尾随的上下文切换将会完成这一工作}、上下文切换\footnote{如上下文切换钩子RMP\_SCHED\_HOOK}或新线程加入延时队列\footnote{如延时插入钩子RMP\_DLY\_HOOK}后都需要检查当前线程的剩余时间，再检查最近超时的线程延时的剩余时间，然后选择两者中较小的一个设置成硬件定时器的下一次超时时间。

内核提供了函数\_RMP\_Tim\_Future用于得到该超时时间。有两种可能的特殊情况需要处理：（1）该超时时间为0，此时说明有线程延时超时，需要立刻手动触发一个定时器中断，（2）如果这段时间的长度超过了定时器能支持的范围，只要将定时器的超时时间设置成允许的最大值即可；RMP的设计保证只有实际经过的时间才会从线程剩余时间中抵扣掉，因此不用担心线程剩余时间丢失。

\begin{apidesc}{\_RMP\_Tim\_Future}{得到最近超时时间片数}
    \apisig{rmp\_ptr\_t \_RMP\_Tim\_Future(void)}
    得到距离未来最近一次超时的时间片数，可根据此设置系统定时器的下一次超时时间。
    \begin{apiparams}
        \apiparam{无。}{}
    \end{apiparams}
    \begin{apiret}
        \apiretval{rmp\_ptr\_t}{距离未来最近一次超时的时间片数。}
    \end{apiret}
\end{apidesc}

在\_RMP\_Tim\_Future计算最近超时的线程延时的剩余时间前，需要保证时间戳RMP\_Timestamp是最新的。否则，相应的计算就会出错。内核为此提供了函数\_RMP\_Tim\_Elapse用于更新该时间戳，在定时器中断外\footnote{指RMP\_SCHED\_HOOK和RMP\_DLY\_HOOK，不包括RMP\_TIM\_HOOK}若要调用\_RMP\_Tim\_Future则必须先调用它。定时器中断中不需要调用它是因为定时器中断本身就会更新该时间戳。

注：观察\_RMP\_Tim\_Elapse调用时点可以发现，它要么出现在上下文切换向量中，要么出现在后续必然进行上下文切换的\_RMP\_Dly\_Ins中，而RMP中线程的时间片总是会在上下文切换后自动补满\footnote{在RME中的时间片不会自动补满，而是需要用户态调度器手动传递时间片}。因此，在\_RMP\_Tim\_Elapse中无需更新当前线程的剩余时间片数量。这进一步导致在上下文切换中，\_RMP\_Tim\_Elapse只要求在\_RMP\_Tim\_Future之前调用，而不需要固定在切换之前调用\footnote{在RME中，同名函数要求必须在向量开头调用，否则一旦切换完成，消耗的时间片就会被归咎给其它线程}。

\begin{apidesc}{\_RMP\_Tim\_Elapse}{记录时间流逝数}
    \apisig{void \_RMP\_Tim\_Elapse(rmp\_ptr\_t Slice)}
    在计算定时器超时时间前记录当前线程已使用的时间片数。
    \begin{apiparams}
        \apiparam{rmp\_ptr\_t Slice}{距离上次填充定时器以来经历的时间片数，该时间片数可通过读取定时器的计数寄存器，并计算其与设置初值之间的差值得到。这个数字不受RMP\_SLICE\_MAX限制，但它不能超过2\textsuperscript{N-1}-1（也即RMP\_MASK\_INTMAX），其中N是RMP\_WORD\_ORDER。上述读取可以先停止定时器再行完成，因为接下来必然会在RMP\_SCHED\_HOOK中发生一次定时器重编程，定时器是否继续运行在此次读取后就无所谓了；RMP\_DLY\_HOOK在调用本函数后还必须将定时器初值和计数值清零，以免后续RMP\_SCHED\_HOOK含有的本函数发生重复的时间抵扣。}
    \end{apiparams}
    \begin{apiret}
        \apiretval{无。}{}
    \end{apiret}
\end{apidesc}

为了进一步节约能源，当系统中只有初始空闲线程在运行，且当前没有正在延时的线程时，就连系统定时器的时钟源都可以关闭。系统特意提供了\_RMP\_Tim\_Idle来探知这一条件。要使用这一功能，需要保证先关闭中断，再使用该函数探测是否可关闭系统定时器时钟源，最后原子性地开启中断并同时进入合适的低功耗状态\footnote{如果CPU不提供原子性开启中断并关闭定时器时钟源的操作，或者中断来临时无法自动令时钟源起振，则此功能无法使用}。具体案例请参见MSP430的移植下的CCS-MSP430G2553工程。

\begin{apidesc}{\_RMP\_Tim\_Idle}{查询系统定时器空闲状态}
    \apisig{rmp\_ret\_t \_RMP\_Tim\_Idle(void)}
    判断是否可以彻底系统关闭定时器及其时钟源以节约能源。
    \begin{apiparams}
        \apiparam{无。}{}
    \end{apiparams}
    \begin{apiret}
        \apiretval{rmp\_ret\_t}{若返回0，可以关闭定时器及其时钟源，若返回-1则不可以。}
    \end{apiret}
\end{apidesc}

值得注意的是，虽然无节拍内核的CPU功耗更低，但其实时性未必更好，因为它势必包含在上下文切换前的时间流逝登记，以及在上下文切换后对定时器的重编程。这都可能增加上下文切换临界区的长度，并对系统的最坏执行时间产生负面影响。在虚拟化环境中，这些缺点更是会成倍放大，因此在非低功耗环境下推荐选择有节拍内核。

\subsection{线程上下文的保存和恢复}\label{subsec:saving-and-restoring-thread-context}

\protect\phantomsection\label{subsec:saving-and-restoring-thread-context-2}在通常情况下，本系统使用一个专用的线程上下文切换中断向量，并在且仅在其中进行线程的上下文切换。这使得上下文切换中断向量成为系统中唯一需要以汇编写成的中断向量，大大方便了系统的移植。这还降低了其它中断向量的中断延迟，因为在进入和退出其它中断向量时不需要压栈和弹栈全部的上下文。此外，在多个中断连续咬尾执行时，上下文的保存、恢复与切换仅在最后一个中断退出后才进行，这节约了很多开销。基于上述原因，推荐使用这种方法。

然而，这种方法并非没有缺点。它（1）需要一个额外的、软件可触发的、可悬起的中断向量\footnote{如ARMv7-M的PendSV}，（2）要求该向量必须在软件触发后几乎立刻进入\footnote{触发后立刻进入中断，或触发到进入中断之间的指令数可控}，（3）要求中断之间可以直接咬尾，也即上一中断退出后可以无缝接续\footnote{部分CPU在中断返回后必须要执行一条主程序指令才能响应下一个中断}下一个，中间不得插入任何主程序指令，（4）增加了从中断向量激活某个线程的延迟，（5）增加了从线程发起上下文切换的延迟。在那些不满足（1）-（3）或不希望（4）与（5）发生的场合\footnote{如将线程当成合作运行的纤程或协程使用时，或者软中断的中断延迟很高时}，可以参考下面的实现思路。

首先，将RMP\_YIELD使用汇编语言编写，并在其中完成如下操作：

\begin{enumerate}
\def\labelenumi{\arabic{enumi}.}
\item
  关中断；
\item
  模拟处理器进入中断的行为，压栈中断会压栈的寄存器，且确保压栈的版本中中断开启；
\item
  将CPU的基本寄存器全部压入线程栈\footnote{指令指针PC（或IP）需要一些特殊处置，压栈时不能压入它的原值，而要压入步骤10的返回点}；
\item
  将现在的线程堆栈指针存入变量RMP\_SP\_Cur；
\item
  切换到内核栈\footnote{这一步是可选的，但推荐使用独立的内核栈和线程栈}；
\item
  调用\_RMP\_Run\_High进行线程切换；
\item
  将变量RMP\_SP\_Cur赋给现在的线程堆栈指针，切换到线程栈；
\item
  将CPU的基本寄存器全部从线程栈弹出；
\item
  调用中断返回指令，开中断并返回；
\item
  调用普通返回指令，返回原程序流执行。
\end{enumerate}

然后，用汇编语言编写中断向量的头尾，并在其入口都加上寄存器的保存（上记步骤3、4、5），出口都加上寄存器的恢复（上记步骤7、8、9），并在出口执行之前检测是否置位。如果置位，执行线程切换（上记步骤6）。宏则无需定义。

使用本方法时，需要保证RMP\_YIELD产生的上下文可以被中断向量的出口恢复，且中断向量的入口产生的上下文也可以被RMP\_YIELD恢复，尤其是在那些进入中断时会进行硬件压栈的架构上。此外，某些架构具备一些执行时间较长的指令\footnote{如ARMv7-M的LDM、STM和IT块，以及部分DSP处理器的硬件循环指令等}，处理器的设计者为了减少中断延迟，允许中断打断它们并记录它们的执行情况以在中断退出时正确恢复它们；对于这类架构，在恢复寄存器时必须能够从它们被打断的地方正确恢复它们，而非从头执行整条指令，否则会导致程序逻辑故障\footnote{如重复压栈、条件判断错误等}。如果不能保证从中间恢复执行却又需要使用本方法，可以在锁调度器后利用提前查看下一线程的上下文，如果该上下文处于被打断的指令中间而无法被恢复，则进入速度较慢的中断上下文切换路线\footnote{参见的移植}。

对于不提供专用上下文切换中断的架构，的官方移植均采取这种方法，并且会将上述汇编程序段封装成宏。所有的系统中断\footnote{会调用\_ISR系列函数的、会访问操作系统临界区的中断}均必须使用该宏来保存和恢复上下文，而透明中断\footnote{不会调用\_ISR系列函数的、与操作系统无关的中断}则无此要求。具体例子请参见的移植。

这一方法还可以被如下各小节所列的措施进一步优化。的官方移植均不实现这些优化措施，因为在综合考虑下它们显得得不偿失。

\subsubsection{容忍不立刻进入的软中断}\label{subsubsec:tolerating-a-software-interrupt-that-does-not-enter-immediately}

虽然我们可以在YIELD调用后都添加对清零的循环等待，并待其清零后才继续执行代码以规避（2）中对几乎立刻进入的需求，但其他需求仍然无法简单规避。因此，在不满足（1）-（3）的架构的官方移植中不采取这个做法。

\subsubsection{推迟关中断以提高实时性}\label{subsubsec:deferring-interrupt-disabling-to-improve-real-time-behaviour}

\protect\phantomsection\label{subsubsec:deferring-interrupt-disabling-to-improve-real-time-behaviour-2}虽然关中断可以推迟到压栈完成后才执行以降低总关中断时间，但若在RMP\_Thd\_Yield引起的RMP\_YIELD的压栈完成后\footnote{由RMP\_Sched\_Unlock引起的RMP\_YIELD调用无此问题，因为此时调度器已锁、中断已关}中断立刻来临，则线程的上下文会被保存两次。类似地，开中断可以提前到弹栈之前就执行，但这样做比推迟关中断更加可怕：打断RMP\_YIELD弹栈过程的中断向量会再次压栈，且此次压栈在未来可能仍由RMP\_YIELD负责弹栈；然而，RMP\_YIELD可能会在来得及弹栈之前再次被中断打断，从而造成理论上无限的栈用量\footnote{在实践中几乎不可能，因此在不严肃的场合无需理会}。为了避免上述问题，RMP的官方移植总会先关中断再进行寄存器压栈，且在弹栈完成之前不会开启中断。

\subsubsection{利用调用约定优化中断延迟}\label{subsubsec:using-the-calling-convention-to-optimise-interrupt-latency}

\protect\phantomsection\label{subsubsec:using-the-calling-convention-to-optimise-interrupt-latency-2}考虑到中断向量未必会在每次被触发时都请求上下文切换，因此在中断向量的入口与出口激进地保存和恢复一切寄存器是不必要的。有两个可能的优化：（1）在入口仅保存那些函数调用约定中调用者保存（Caller-Save）的寄存器，然后立刻执行中断向量，待执行完毕后在向量末尾根据RMP\_Sched\_Pend判断是否需要切换上下文。如果确需上下文切换，则再压栈剩余的寄存器并切换，否则则直接弹栈调用者保存的寄存器并退出中断。（2）进一步地，在向量末尾判断退出后是否还有其它悬起的中断，若是则可将上述判断推迟到最后一个向量退出时才执行，本次无论RMP\_Sched\_Pend是何值都不执行上下文切换\footnote{以达到类似ARMv7-M的PendSV的效果}。但是，这两个优化（1）需要适应各种编译器的潜在的不同的调用约定，（2）增加硬件抽象层的复杂程度，（3）仅提高平均效率而恶化最坏执行时间，因此在RMP的官方移植中不予实现。

\subsubsection{利用调用约定优化出让效率}\label{subsubsec:using-the-calling-convention-to-optimise-yield-efficiency}

由于是一个实为汇编函数的宏，编译器在调用它前必然遵循函数调用约定，以将调用者保存的寄存器提前保存起来，并在函数退出后自动加以恢复。事实上，仅需要保存和恢复被调用者负责保存（Callee-Save）的寄存器；这可以提高线程间自愿出让处理器的效率，类似的设计在宏内核中相当常见。然而，RMP\_YIELD产生的上下文将不再完整，中断向量处的上下文保存与恢复必须想办法与它兼容。不管使用什么兼容手段，中断延迟都将毫无疑问地恶化，因而此优化在的官方移植中不予实现。

注1：一个讨巧的兼容手段是让中断向量使用线程栈，并且变造需要切换上下文时的中断返回路径使其（1）提前让处理器中断复位\footnote{如ARMv7-M等会记录中断状态且自动发送中断结束信号（End-Of-Interrupt，EOI）的处理器需要}但仍保持中断关闭，（2）定制一个末状态关中断\footnote{它原本不应该在中断向量中调用，且在末尾总会开启中断}的变种并调用之，以生成一个可被正常RMP\_YIELD恢复的不完整上下文，最后（3）弹栈退出中断。再结合\ref{subsubsec:using-the-calling-convention-to-optimise-interrupt-latency}
的优化\footnote{当然也可以不结合，这会使中断压和弹完整的栈，无谓地恶化性能}，使中断入口和出口仅负责压栈和弹栈调用者保存的寄存器，一切被调用者保存的寄存器的保存与恢复都由负责，绕开了兼容性问题。然而，这样做的代价是（1）必须为每个线程栈预留中断向量用量，增加内存消耗，和（2）下次切换回来时需要先弹不完整的栈再弹中断栈，多了几次跳转，恶化最坏执行时间。

注2：在上述兼容手段中，某些实现为了最小化关中断时间走得更远。它们会在（1）处理器中断复位后就放任其处于开中断状态以便在被调用前响应累积的悬起中断，而且会在（2）处直接调用原版，这样就能在中断出口的弹栈执行前就响应中断。这的确减少了关中断时间，但可能会导致来得及执行前或者刚刚执行结束后该线程又被中断打断压入寄存器组，造成类似\ref{subsubsec:deferring-interrupt-disabling-to-improve-real-time-behaviour}
的理论上无限的栈用量。如果还同时施加了\ref{subsubsec:deferring-interrupt-disabling-to-improve-real-time-behaviour}
的优化，则栈用量情况还将进一步恶化。

\subsubsection{出让与抢占间的代码复用}\label{subsubsec:reusing-code-between-yield-and-preemption}

在的官方移植中，IELD与中断向量的出入口总是使用两套独立的实现，无论它们有多么相似也不复用任何代码。这是因为（1）我们不能假设所有汇编器都允许包含汇编头文件；在允许的场合，我们提供一个可供用户包含到中断向量汇编文件中的文件，该文件也可被架构的底层汇编文件包含，此时可以实现复用；在不允许的场合，提供一个可供用户手动粘贴到中断向量汇编文件中的文件，此时无法实现复用；（2）这两套实现在某些架构上大相径庭，无法复用，比如某些架构在进入和退出中断时会由硬件执行栈对齐等特殊操作\footnote{如ARMv7-的移植等}，而IELD必须用另辟蹊径来模拟这套操作。总的来说，与其仅针对部分架构和工具链使用一套复用实现，不如对所有架构都使用两套独立实现；这能够保证所有底层心智模型的整齐划一。

\subsection{协处理器上下文的保存和恢复}\label{subsec:saving-and-restoring-coprocessor-context}

协处理器指浮点处理器、向量处理器等CPU扩展模块，它们也有自己的上下文；如果多个线程要并发共享同一个协处理器，则必须在线程上下文切换时一并保存和恢复该协处理器的上下文。为此，提供了TXAVE和TXAD两个钩子。在保存钩子内，先要探测当前线程是否使用了协处理器，如果是则保存其上下文；在恢复钩子内也要先探测协处理器是否被使用，如果是的话则恢复其上下文。如果无法探测该协处理器是否被使用，则必须总是保存和恢复其上下文。

由于这两个钩子会在进入内核的语言环境后被调用，在具备独立内核栈的移植中可以直接利用来将上下文保存在线程栈上；的所有官方移植都保证使用独立的内核栈。在内核共用线程栈的移植上则必须将协处理器上下文保存在额外的结构体中，否则可能会覆写内核返回地址而导致出错。

对于绝大多数架构，的官方移植中已经在栈上实现了其协寄存器的上下文保存和恢复。对于这一部分架构，不需要使用钩子函数就可以保存和恢复其浮点寄存器的上下文。针对不同的浮点处理器配置，可能需要在汇编文件中引用不同的保存和恢复例程。

注1：在的官方移植中，一切通用寄存器在内的上下文都被放在线程栈上。这是因为（1）几乎所有的处理器在进入中断时都至少会将和压栈\footnote{仅有的例外是和}，部分处理器还会强制压栈一部分通用寄存器\footnote{如系列、等}，如果要将一切上下文都放进结构体中则需要拷贝这些寄存器到结构体，于性能不利，且（2）若将未自动压栈的上下文单独保存进结构体中，则线程启动前既需要需初始化线程栈中所含的和，又需要初始化结构体，于软件工程不利。

注2：在的官方移植中，为了使所有浮点处理器配置情况都能收录收进一个汇编文件中，对浮点处理器的操作一律采用机器码。这可以在浮点处理器不存在时规避汇编器报错；会根据配置情况引用正确的保存和恢复例程，因此不会引起未定义指令异常。

注3：在RMP上，推荐的协处理器使用方法是仅让一个线程使用它，这样就可以不必保存和恢复它的上下文，从而提高系统的效率和实时性。对于极少数特殊架构\footnote{如DSPIC系列的指令相关寄存器}，其协处理器的上下文由硬件全权控制，无法被软件保存和恢复，此时请务必保证系统中只有一个线程使用该协处理器。

\subsection{线程内存保护的实现}\label{subsec:implementation-of-thread-memory-protection}

在微控制器中，通常利用内存保护单元（Memory Protection
Unit，MPU）来实现内存保护。与能进行地址翻译的内存管理单元（Memory
Management
Unit，MMU）不同，MPU仅能限制物理内存的访问权限，无法进行地址翻译。MPU由一系列的寄存器组成，这些寄存器里面填充着各个被保护的物理内存区域的访问属性。每个线程可以拥有一套不同的寄存器设置，并以此来限制自身对内存的访问。

推荐的实现是将MPU看作协处理器，也在TXAVE和XAD内进行其上下文的保存和恢复。需要注意的是，每个线程的MPU设置都应当允许读写的内核内存和执行的代码；这是因为没有真正的用户态和内核态，而且所有的系统调用都被直接实现为函数调用。正因如此，基于实现的内存保护仅能增强系统的功能安全性\footnote{又叫可靠性；如探测栈溢出等常见事故}，但不能提供任何信息安全性。

\subsection{低功耗设计注意事项}\label{subsec:low-power-design-considerations}

在低功耗设计中，推荐使用无节拍内核，此外还要在钩子内插入能让处理器进入休眠状态的指令\footnote{在ARM上，WFI和指令可以达到这个效果，而在上则必须修改寄存器的某些位}。

\section{RMP中已知的影响实时性的因素}\label{sec:known-factors-that-affect-real-time-performance-in-rmp}

在中已知的影响实时性的因素只有两个，分列如下。这两个因素在实际的应用程序中都很好避免；通常而言，如果因为这些因素使系统的实时性受到损伤，应用程序就不是良好设计的。

\subsection{延时队列}\label{subsec:the-delay-queue}

使用一个延时队列管理所有的延时。这个延时队列从前向后，各线程的超时时间依次增加。这样，每次时钟中断，都只需要检查队列中的第一个线程即可。但是，这个设计要求每当插入线程时都要遍历队列找到合适的插入位置以使得线程超时时间仍然递增。这个遍历的最坏时间复杂度对队列中已有的线程数量是O(n)的\footnote{有数据结构存在，但的应用场景一般不需要如此重量级的实现}。从数学上可以证明，不存在一种数据结构能够使得定时器的所有操作总为O(1)，因此这必然要求应用程序不要特别大量地使用延时等待或延时阻塞。

此外，提供的延时总是系统级的，低优先级线程的延时操作也需要锁调度器进行处理，这可能会干扰高优先级线程的运作。如果这会造成问题，更推荐用户使用有限的几个线程实现软定时器功能，并分别服务不同优先级线程的延时。

\subsection{删除内核对象}\label{subsec:deleting-kernel-objects}

线程邮箱和信号量等部分内核对象可以拥有一个等待队列，而删除该内核对象时需要将该等待队列中的线程从该队列中一一去除。如果某个内核对象上有过多线程在等待，则删除内核对象可能导致过久的锁调度器时间。

\subsection{锁调度器和关中断}\label{subsec:locking-the-scheduler-and-disabling-interrupts}

提供了锁调度器和关中断的接口。但是，这对系统实时性是有显而易见的损害的，因此在应用程序中应当尽量避免此类操作。比如，如果已知只有一个线程会在某内存池中做操作，那么再分配和释放内存时就可以不用锁调度器。

\section{降低RMP的中断延迟}\label{sec:reducing-the-interrupt-latency-of-rmp}

为了保证临界区不被入侵，在进行特定操作时需要锁调度器，而它的最简单的实现是关闭总中断。然而，这样做会使整个系统失去对外界事件刺激的即时响应，即使是那些与系统无关的透明中断\footnote{不会调用\_ISR系列函数的、与操作系统无关的中断}也会遭到禁用，导致从刺激源发生到中断向量执行之间的延迟变长\footnote{这里的中断延迟指的是硬件中断延迟，它和RTOS的中断延迟是两回事；后者一般是指从中断中进行发送到对应的接收线程被激活的延迟，是RTOS的固有值，无法被降低}。在实时性要求较高的机电设备或接口模拟等应用场合，这种行为是不允许的。为此，提供了一系列方法来降低或彻底消除由临界区导致的中断延迟。无论使用哪种方法，系统中断\footnote{会调用\_ISR系列函数的、会访问操作系统临界区的中断，包括时钟中断和专用上下文切换中断}都必须使用RMP提供的上下文保存宏，而透明中断则只要求保存和恢复任何会被它们修改的寄存器，且系统中断的抢占优先级必须被设为相同以使其不能互相嵌套。

\subsection{按优先级掩蔽系统中断}\label{subsec:masking-system-interrupts-by-priority}

如果架构\footnote{如Cortex-M系列等}提供了掩蔽某优先级以下中断的功能，且系统中断的优先级都在某值以下，就可以通过掩蔽优先级在该值以下的中断来达到仅禁止系统中断而不禁止透明中断的效果。在这些架构，RMP的官方移植均使用掩蔽来替代禁止总中断，天然地不会引入任何中断延迟。配置头文件中的掩蔽等级则可供用户视具体情况调整。

\subsection{按中断源掩蔽系统中断}\label{subsec:masking-system-interrupts-by-source}

\protect\phantomsection\label{subsec:masking-system-interrupts-by-source-2}在那些无法按优先级掩蔽中断的架构上，也可以通过依次禁止各系统中断源来达成和按优先级掩蔽相同的效果。不过，由于无法预测用户要使用哪些系统中断源，RMP无法提供统一的方案：官方移植中的RMP\_INT\_MASK()和RMP\_INT\_UNMASK()宏都会操作总中断，RMP\_YIELD也会在进入和退出时操作总中断。

为了达到按中断源掩蔽中断的效果，需要对官方移植做如下修改：

\begin{enumerate}
\def\labelenumi{\arabic{enumi}.}
\item
  将RMP\_INT\_MASK()和RMP\_INT\_UNMASK()宏分别指定为对各系统中断源的屏蔽操作和解除屏蔽操作。
\item
  在进入中断后，保存中断返回地址\footnote{绝大部分架构在进入中断时会自动将返回地址压栈，如等则需要手动压栈后再开启中断，否则透明中断可以在返回地址还没来得及保存之前就覆写它而导致系统中断无法正确返回其原有上下文}，并立刻开启总中断。这只需要在调用RMP的官方上下文保存宏之前插入一条开启总中断的指令即可，这样关闭总中断的时间最多只有一条指令。如果硬件不提供将所有系统中断都设置为相同的抢占优先级的功能，则需要对提供的宏进行一定修改，以在开启总中断之前屏蔽各系统中断源，并在中断返回之前（1）禁用总中断，（2）解除对各系统中断源的屏蔽，以及（3）原子性地使能总中断并从中断中返回\footnote{中断返回指令可以做到}。
\item
  在IELD的开头和结尾加上对系统中断源的屏蔽操作和解除屏蔽操作。
\end{enumerate}

\subsection{委托操作给软件中断}\label{subsec:delegating-operations-to-a-software-interrupt}

\protect\phantomsection\label{subsec:delegating-operations-to-a-software-interrupt-2}在系统中断源较多或者中断控制器配置不便的场合，对各系统中断源一一屏蔽和解除屏蔽会造成巨大的移植和执行负担。此时，如果处理器具备可悬起的软件中断源，或者软件可触发的多余的外设中断源\footnote{外设的中断源一般总会富余，此时可以拿来做此用途}，则更高优先级的中断可以选择将部分操作委托（Delegate）到它内部执行，此时该软件中断是系统中断，而更高优先级的中断则是透明中断，大大降低了系统中断源的数量。如果中断发生得较频繁，这种做法还有降低处理器负担的优点，因为透明中断只需要保存和恢复它们使用的寄存器\footnote{在那些具备多组寄存器的架构上，透明中断可以使用其它寄存器组以彻底绕过保存和恢复}。不过，对于被委托的操作而言，其本体的中断延迟时间和到线程的中断延迟都变长了，因此本方法有利有弊。

为达到委托操作的效果，需要对官方移植做如下修改：

\begin{enumerate}
\def\labelenumi{\arabic{enumi}.}
\item
  按照\ref{subsec:masking-system-interrupts-by-source}
  中的方法，将包括该软件中断源在内的一切系统中断添加进按中断源掩蔽的对象中。
\item
  在透明中断中，向某数据结构记录要进行的系统操作，并触发该软件中断源。
\item
  在该软件中断中，从某数据结构读取系统操作，并代替透明中断完成它。

  在该方法中，数据结构的设计是一个难点，因为该数据结构可能会被透明中断和软件中断同时修改。如果数据结构是较为简单的计数器等，可以使软件中断保留一个本地副本，并计算本地副本和计数器当前值的差值来决定应当进行多少次操作；计数器的原子性读取则可通过循环读取直至稳定的方法实现。如果数据结构需要记录复杂的额外信息，则可以采用原子操作指令访问它，甚至可以使用无锁（Lock-free）或无等待（Wait-free）数据结构；在处理器不提供这种操作时，也可以通过短暂地禁用对应的透明中断来达成相同的目的。有关互斥方案的设计请自行参阅相关资料，在此不做过多介绍。
\end{enumerate}

\subsection{委托操作给定时器中断}\label{subsec:delegating-operations-to-the-timer-interrupt}

在找不到或不方便使用软件可悬起中断源又希望进行操作委托的场合，还可以选择将操作委托给定时器中断。它的原理和委托给独立的软件中断是类似的，只是（1）定时器中断负责处理委托，（2）定时器中断不由透明中断触发而由定时器硬件触发，且（3）系统必须是有节拍内核。

在实现定时器委托时，需要对官方移植做如下修改：

\begin{enumerate}
\def\labelenumi{\arabic{enumi}.}
\item
  按照\ref{subsec:delegating-operations-to-a-software-interrupt}
  中的方法处理定时器中断这个唯一的系统中断源。
\item
  在透明中断中，向某数据结构记录要进行的系统操作。
\item
  在系统嘀嗒钩子IMOK中，从某数据结构读取系统操作，并代替透明中断完成它。
\end{enumerate}

此方法相当于将每个时间片内的所有中断操作都推卸给了定时器中断，因此这些操作可能被延迟很久；只有当线程并不需要在立即获知中断发生时才能使用这个方法。如果这些操作是通知性质而非计次性质，则可将多次操作合并为一次操作，此时能获得可观的性能提升\footnote{尤其是针对缓冲区较深的网卡、串口等设备，可以等到接收到一定数据量时再行批次处理}。

\subsection{限制不关键中断对关键中断的干扰}\label{subsec:limiting-interference-from-non-critical-interrupts-on-critical-interrupts}

在混合关键度系统中，部分透明中断源的关键性不高但优先级却很高\footnote{紧急而不重要，偶尔丢失中断无所谓，如遥控解码、流媒体传输等}，如果不对其中断频率加以限制，必然会消耗大量的CPU，甚至干扰关键但相对不紧急的其它中断的执行，造成它们的中断延迟的不受控的增长。

为了限制这些透明中断源的频率，需要采取以下措施：

\begin{enumerate}
\def\labelenumi{\arabic{enumi}.}
\item
  为每个透明中断源准备一个计数器，每发生一次中断则原子性地将其自减一，当计数器递减至零时自行关闭中断源。
\item
  每当定时器中断或关键中断发生时，原子性地为该计数器（1）赋恰当的初值或者（2）在其未达上限之前增加一个定值，并打开对应的中断源。

  此时，每个透明中断源都相当于被限制在了一个可推迟服务器（Deferrable
  Server）或偶发服务器（Sporadic
  Server）内，在一定的时间段内只能被触发有限的次数，如果有多于这个次数的触发，则中断会丢失而不予响应。这样一来，低优先级但高关键度的中断所受的干扰就是可控的。
\end{enumerate}

\section{缩减RMP的存储器占用}\label{sec:reducing-the-memory-footprint-of-rmp}

在内存特别宝贵的场合，有必要小心配置以节省本就不多的RAM和ROM。此种状况多见于等16位单片机和等低端32位单片机上。在实践中，可以采用如下手法缩减的内存占用：

\subsection{降低抢占优先级的数量}\label{subsec:reducing-the-number-of-preemption-priorities}

不同于其他的操作系统，的线程邮箱是内建于线程之中的零成本功能，而其它功能在不使用时则不占用任何存储。因此，的编译时选项中只有一个会决定其内核的RAM占用，那就是支持的优先级个数\footnote{也即}。由于的每个优先级都使用一个双向循环链表来管理，因此每个优先级都会消耗2个机器字长的RAM。在一些微型系统上，配置过多的优先级是完全没有必要的，因此可以将支持的优先级调整成3个来最大限度节省RAM。在这种用例下，本身的RAM消耗甚至可以被挤压到128字节以内\footnote{如果不算内核中断响应用堆栈，那么实际大小是64字节}。

\subsection{调整厂商提供的库}\label{subsec:adjusting-the-vendor-provided-library}

某些厂商提供的HAL库和启动文件可能会静态分配一些堆和栈\footnote{比如Cortex-M系列微控制器的汇编启动文件的最前面}。在使用RMP的场合，厂商分配的堆是完全没有必要的，可以置0；而厂商分配的栈则会被OS内核在响应中断时使用\footnote{也即作为内核栈使用}，因此只要确认中断处理程序的最大栈消耗，然后将厂商分配的栈大小设置到比这一数字稍多\footnote{一般推荐多几十字节}即可。此外，厂商库中用不到的文件可以不添加进工程，这也有助于减少RAM和ROM用量。如果厂商库占用的存储器不可忍受，那么也可以选择直接操作外设寄存器。

\subsection{调整编译器选项}\label{subsec:adjusting-compiler-options}

通过调整编译器选项，也有可能减小ROM和RAM占用。调整编译器选项为代码体积优化对减小ROM占用尤其有效，对减小RAM占用有时也有效果。

\subsection{不使用动态内存管理}\label{subsec:not-using-dynamic-memory-management}

内建了一个动态内存分配器。该分配器在小内存设备上不推荐使用，因为分配器本身也会有一定的固定内存消耗。在小内存设备上，推荐静态分配所有的内核对象和变量。

\section{移植时的推荐规则}\label{sec:recommended-rules-for-porting}

在部分架构上，移植RMP可能会遇到两难抉择。下文将介绍这些抉择为何困难并给出推荐的处理规则。

\subsection{处理器字长的推荐规则}\label{subsec:recommended-rules-for-processor-word-length}

在32位或更高的处理器上，指针的长度一般都是一个机器字长，这是因为这些处理器普遍能够以任何寄存器做间接寻址，而且机器字长的长度能够寻址整个地址空间，此时可以直接将rmp\_ptr\_t定义为一个机器字长。然而，对于字长更短的处理器，它们的（1）数据与代码空间可能是分开的，（2）指令集设计限制很大，（3）处理器字长不足以寻址整个地址空间，因此决定rmp\_ptr\_t的长度需要一些技巧。

\subsubsection{8位处理器的情况}\label{subsubsec:the-case-of-8-bit-processors}

在绝大多数8位处理器上无法或不需要移植RMP，这是因为它们（1）内存资源太少，（2）为追求内存利用效率\footnote{也有可能是由于架构缺乏SP相对寻址功能}，编译器采用局部变量静态覆盖技术，（3）为追求硅片面积效率，函数调用栈是由硬件实现的，（4）应用程序较为简单，使用RMS状态机就可以表达。在少数资源丰富且可以移植RMP的8位处理器上，几乎所有的编译器都默认指针长度为16位。这是因为8位仅能寻址256字节，这对任何应用程序来说都太少了。因此，rmp\_ptr\_t至少要定义为16位。

极少数8位处理器甚至允许寻址24位内存空间，此时需要分成两种情况讨论。如果数据指针或函数指针中任意一个的内部二进制表示\footnote{RL78等部分架构的指针字面值是20位，但内部二进制表示却是16位}高于16位，则在该架构上rmp\_ptr\_t需定义为32位，这虽然会浪费空间但别无他法。如果两个指针的内部表示都是16位，则在该架构上rmp\_ptr\_t可定义为16位。后一种情况是可能的，因为高于16位内存空间的寻址有两种实现方法：（1）由独立于指针的段寄存器实现，而若编译器的指针不包括段寄存器则指针就是16位\footnote{如DSPIC等架构}，且段寄存器的切换需要用户程序手动进行，（2）工具链保证将内核使用的变量和代码链接到前64KiB地址空间\footnote{如AVR等架构}，后面的地址空间则可供绝大多数不被指针引用的变量和代码使用。有关段寄存器和内存模型的详细内容，请参见下两个小节。

\subsubsection{16位处理器的情况}\label{subsubsec:the-case-of-16-bit-processors}

绝大多数16位处理器都具备移植RMP的可能。这些处理器也可以分为只能寻址16位地址空间的类型，和能寻址更多地址空间的类型。对于前者，将rmp\_ptr\_t定义为16位即可；对于后者，则可视高内存空间的寻址方法，仿照8位处理器的方案处理。

\subsection{段寄存器的保存和恢复规则}\label{subsec:rules-for-saving-and-restoring-segment-registers}

部分处理器具备段寄存器\footnote{如等架构}，部分编译器也可能在某些条件下使用一个全局变量指针寄存器\footnote{如、等架构}。根据编译器对其使用风格的不同，段寄存器的上下文保存和恢复也可做不同处理。

\subsubsection{段寄存器仅被启动代码设置}\label{subsubsec:the-segment-register-is-set-only-by-startup-code}

这是最容易应对的情况。此时，段寄存器仅在启动时被启动代码设置一次，之后便不再进行更改；这可能是由于程序较小，在一个段内即可放下。如果能确信程序的行为确实如此，则可以完全不进行段寄存器的保存和恢复。

\subsubsection{段寄存器总在访存前临时更改}\label{subsubsec:the-segment-register-is-changed-temporarily-before-each-memory-access}

这是较容易应对的情况。此时，编译器不会对段寄存器的值做任何假设，而是在访存前临时将段寄存器指向需要的段。由于段寄存器的值可能发生变动，且中断可能插在段寄存器修改和段寄存器使用之间，因此需要对段寄存器做上下文保存和恢复。

\subsubsection{段寄存器在函数调用前要求保持初值}\label{subsubsec:the-segment-register-must-hold-its-initial-value-before-a-function-call}

这是最难应对的情况。此时，编译器不仅会在函数内部生成修改段寄存器的操作，还假设在进入包括中断向量在内的每一个函数之前段寄存器都是初始值；这一般是为了优化程序的效率，若启动代码能将段寄存器在初始化时就指向访问最频繁的段，则可在函数调用之间省去段寄存器设置工作，并减小代码体积。那些访问其它段的函数则需要在它们的头尾保存和恢复段寄存器。

在此种情况下，不仅需要对段寄存器做上下文保存和恢复，还需要在中保存段寄存器的初值；在中断向量\footnote{和IELD中，如果不使用专用线程上下文切换中断向量的话}中如果需要调用函数，则必须在调用前将各段寄存器设置为保存下来的值。这种做法的优点在于，不论工具链和启动代码将这些值设置成什么，它们都必然被截获，这样就不需要从五花八门的链接器脚本中取得这些值。当然，如果工具链是等开源标准品，也可以选择从链接器脚本中直接取得这些值。

部分工具链的段寄存器使用风格随着选择的内存模型而变化。因此，的几乎所有官方移植都采取上述方法以消除段错误风险，即便这可能是多此一举。只有一种例外：该架构只计划提供一种工具链支持，且该工具链明示了各段寄存器的初值为某固定值\footnote{如仅提供工具链移植的架构}。

\subsection{内存模型的推荐规则}\label{subsec:recommended-rules-for-memory-models}

在部分8位和16位处理器上，由于有段寄存器的存在，程序可以选择使用一个或多个代码段（CODE）、数据段（DATA）或堆栈段（STACK），从而产生不同的内存模型（Memory
Model）。当程序使用一个段时，其指针为16位，但其寻址范围被限制在一个段以内；当程序使用多个段时，其指针为24或32位，其寻址范围包括整个内存地址空间。部分架构甚至将同一种存储器按地址空间分为多种类别，每种类别的指针长度都各不相同。此外，代码段、数据段和堆栈段的默认指针类型是可以独立设置的：假设代码段的指针长度有C种，数据段的指针长度有D种，堆栈段的指针长度有S种，则内存模型的种类有C×D×S种。部分架构的段寄存器方案甚至到了五花八门的地步\footnote{如基于的红白机（Famicom/NES）的存储器管理器（又称Mapper）}，因为它们并未提供架构级的段寄存器，而是要求用户在外部使用分立逻辑配合I/O口实现存储器映射电路。

为了简化起见，RMP不可能针对每一种情况都给出相应的移植，且RMP在设计时假设一切指针均为等长，这要求代码段、数据段和堆栈段的默认内存模型相同。这可能在移植时导致两难选择：如果总是不加考虑地选择最大的内存模型，则rmp\_ptr\_t是32位，大大增加了处理器的负担，尤其是程序规模较小时得不偿失；如果总是不加考虑地选择最小的内存模型，则rmp\_ptr\_t是16位，处理器负担较轻，但无法访问更多地址空间。为此，我们按照以下原则将处理器分为三类，并针对不同的类别分别给出推荐方案。

\subsubsection{常用内存模型为大模型}\label{subsubsec:the-common-memory-model-is-the-large-model}

小部分处理器\footnote{如8051、80251等架构}从诞生之初就常使用大内存模型，它们的绝大多数工具链默认sizeof(void*)为4。这是由于它们的存储器类别太多、部分类别的寻址范围又太少，就连工具链作者也难以应付大量组合带来的复杂性，只能采取长指针来逃避它。

在这种处理器上，的官方移植采取大模型：它将rmp\_ptr\_t定义为32位，其中高16位就是段寄存器的值。段寄存器和相应的裸指针一并保存和恢复。

\subsubsection{常用内存模型为小模型}\label{subsubsec:the-common-memory-model-is-the-small-model}

\protect\phantomsection\label{subsubsec:the-common-memory-model-is-the-small-model-2}大部分处理器仅在用户明示时才使用大内存模型，它们的绝大多数工具链默认sizeof(void*)为2，这是为了尽量提升程序性能敏感部分的效率。当程序需要访问更多内存时，程序员需要在变量和函数前使用far等关键字声明其为远符号，编译器会自动生成相应的段和跨段访问代码，使其在小模型下也能访问更宽广的内存地址。

在这种处理器上，的官方移植采取带段寄存器保存和恢复的小模型：它将定义为16位，但会在上下文切换时切换段寄存器以支持线程对其它段的访问。用户必须确保所有的内核数据结构\footnote{包括各对象控制块以及消息队列中的消息}都与内核的数据段是同一个，且所有的线程入口\footnote{不将它们声明为即可}都与内核的代码段是同一个；不过，一旦线程启动，它可以访问其它段的变量以及跳转到其它段去运行，不受这一限制。

在移植时，堆栈段需要视情况做特殊处理：（1）如果该处理器未提供堆栈段寄存器，或其组成的常见系统中RAM与64KiB相去不远\footnote{常见于微控制器：它们要么没有堆栈段寄存器，要么内存小于64KiB而使堆栈段寄存器总为固定值}，则堆栈段无需特殊处理；（2）如果该处理器提供了段寄存器且其RAM总量常远超64KiB\footnote{常见于等通用微机系统}，此时可以将宏定义为1以将堆栈段寄存器保存在线程控制块中。一旦该宏被启用，``创建线程''系统调用和``线程栈初始化''底层调用的格式均会发生改变，在参数前会出现一个额外的ment参数以方便用户传入堆栈段寄存器的值。

注：在（2）中，堆栈段寄存器不能保存在栈上，因为从栈中恢复堆栈段寄存器需要弹栈，而弹栈又会用到堆栈段寄存器，从而造成悖论问题。

\subsubsection{常用内存模型无法推测}\label{subsubsec:the-common-memory-model-cannot-be-inferred}

一部分处理器的内存模型无法推测\footnote{如、等架构，但不包括等官方定义了段寄存器的架构}，它们的绝大多数工具链虽默认为2，但允许用户编写自定义的段切换代码以配合片外定制电路实现跨段访问。

对于此类架构，的官方移植一律采取不理会段寄存器的小模型。用户必须对底层移植进行少许定制修改以适应其硬件结构。

\subsection{移植兼容性的推荐规则}\label{subsec:recommended-rules-for-porting-compatibility}

几乎所有微控制器厂商都会在一个基础指令集上搭配一系列可选的协处理器\footnote{或者扩展指令集等}以适应不同的细分市场。部分厂商甚至会增加老架构中寄存器的长度\footnote{如MSP430X等架构}来适应新的需求。在移植RMP时，一个简单的办法是针对不同的协处理器组合编写独立的移植，但这样会导致各移植之间包含巨量的重复代码。因此，我们制定一些规则，并依此判断哪些协处理器组合应该被合并为一个移植。

\subsubsection{链接能跑原则}\label{subsubsec:the-link-and-run-rule}

这是决定合并移植的最重要原则。如果多种协处理器组合\footnote{含无协处理器的情况}的二进制中存在一个子集，该子集可以被链接到任何一种组合，且得到的最终二进制能够在该组合上运行，则这些组合可以合并为一个移植。此时，在各芯片的配置头文件中额外增加协处理器选择宏RMP\_XXX\_COP\_YYY，其中XXX是移植名称，YYY是该扩展指令集，选中该宏则选中该协处理器组合。

协处理器之间的相容性可以分为``含''性、``与''性、``或''性、``非''性和``异或''性：对于A和B两个协处理器，``含''性意味着若A存在则B必须存在，反之未必成立；``与''性意味着若A存在则B必须存在，反之亦然；``或''性意味着A和B的存在性相互独立，互不干涉；``非''性意味着A和B中仅能存在一个，但可以都不存在；``异或''性意味着A和B中必须且仅能存在一个。不同组合之间的相容性由移植底层代码在编译时确认，并在用户选中错误的组合时通过\#error报告。通常而言，微控制器的协处理器都很简单，不涉及非常复杂的兼容性判断。

注：``链接能跑''原则除涉及各协处理器组合的指令公共子集关系之外，还涉及到主流编译器对各协处理器组合的处理方式。在以下情况中，即便存在公共的指令子集，软件上也可能互不兼容：（1）该公共子集图灵不完备以至于无法单独形成应用程序，（2）该公共子集无法在编译器中单独选择，也即无法通过编译生成仅含有该公共子集的二进制\footnote{汇编也许可以，但它的编写效率太低，因而毫无意义}，（3）不同协处理器组合的调用约定完全不同以至于无法以统一方式传递两个以内的参数，（4）不同协处理器组合的内存模型完全不同以至于无法兼容，以及（5）不同协处理器组合要求的栈对齐字节数不同。如果出现情况（1）-（4），则认为不满足本原则\footnote{比如C54x和C55x由于出现（3）、（4）的情况，被认为是两个独立的架构}，必须为每种组合提供独立的移植；如果仅出现情况（5）则认为满足本原则，且提供的合并移植中的栈对齐字节数将是所有组合的可能最大值\footnote{比如C64x、C67x、C64x+、C67x+、C66x由于仅出现（5）的情况，被合并为C6x架构}。

\subsubsection{心智统一原则}\label{subsubsec:the-mental-model-rule}

这是决定合并移植的次要原则。如果以下两条中的任意一条原则成立，则认为该协处理器组合具有统一的心智模型，则这些组合可以合并为一个移植：（1）厂商在其官方资料中明示这些协处理器组合是同一种处理器的衍生品\footnote{如ARMv7-M的FPV4\_SP、FPV5\_SP、FPV5\_DP}，（2）厂商虽未明示，但在技术手册中经常将这些协处理器组合以``合订本''的形式相提并论\footnote{如DSPIC的PIC24、PIC33及其他五花八门的衍生品}。

对于满足``链接能跑''原则，但不满足``心智统一''原则的协处理器组合\footnote{如ARMv6-M和ARMv7-M}，RMP的官方移植总是将其分开以维持其心智模型的独立性。

\subsubsection{高度重复原则}\label{subsubsec:the-high-repetition-rule}

这是决定合并移植的最次要原则。如果以上两条均成立，且协处理器组合之间的底层代码有相当程度的重复度，则这些组合可以合并为一个移植：高重复度说明（1）协处理器组合之间的相似度高，（2）这些协处理器组合对应的产品处于相同或相近的市场生态位。

对于满足``链接能跑''和``心智统一''原则，但不满足``高度重复''原则的协处理器组合，的官方移植总会将代码合并在一个文件内，但汇编文件则可能分开提供。此时，在各芯片的配置头文件中也额外增加协处理器选择宏RMP\_XXX\_COP\_YYY，其中XXX是移植名称，YYY是该扩展指令集，选中该宏则选中该协处理器组合。

\section{厂商库的兼容性}\label{sec:vendor-library-compatibility}

在开发者时间极为宝贵的今天，绝大多数微控制器厂商都提供了固件库。通过这些易学易用的库函数访问片上外设要比使用寄存器读写容易得多。不幸的是，部分固件库的底层功能可能会与操作系统冲突，因而需要修改\footnote{部分厂商考虑到了这一点，将默认的实现加上了``弱符号''属性，此时只要重写它们即可}。常见的冲突发生在中断以及延时两方面，其解决方案如下所述。

\subsection{开关中断的兼容性}\label{subsec:interrupt-enable-disable-compatibility}

操作部分外设时，其读写序列必须在规定的时间之内一次性做完，在此期间被中断打断是不允许的。此时，固件库往往会先关中断以便完成操作，然后再开启中断。对于这些库函数，不要在系统锁调度器期间调用它们，否则它们会在调度器解锁之前不通知操作系统而开启中断，造成系统崩溃。如果无法确定这些库函数是否会在锁调度器期间被调用，可考虑将其开关中断函数更换为提供的调度器锁定/解锁函数，并将RMP\_INT\_MASK和RMP\_INT\_UNMASK指定为掩蔽所有中断。

\subsection{延时函数的兼容性}\label{subsec:delay-function-compatibility}

操作部分外设时，其读写序列往往需要等待外设就绪，因此步骤间常插入延时函数。固件库实现延时函数的方式是五花八门的，可以采用忙循环法、定时器查询法或定时器中断法，部分固件库中甚至针对不同的用途同时采取多种实现方式。

忙循环法的延时函数通过运行一段无用的指令序列来达成延时，它们的行为与一般代码无异，不需要做任何修改就能与RMP兼容。这种延时函数还能用于锁调度器或关中断的场合，因为它不依赖于中断系统。当然，如果能够确定固件库从不在锁调度器或关中断期间调用这种延时函数，可以考虑将它们重写，并使用更节约处理器性能的RMP\_Thd\_Delay替换之。

定时器查询法的延时函数通过初始化定时器并反复查询定时器的情况来确定延时是否完成。它比忙循环法更精确，也同样能在锁调度器期间使用，但它们可能干涉RMP的节拍定时器。此时，该延时函数可能会修改节拍定时器的（1）计数器值，（2）计数上限值或者（3）工作频率，这都可能导致系统工作不正常。更有甚者还会在初始化时关闭定时器的中断，导致RMP的调度器失控。当上述情况发生时，可以重写该函数，并（1）使其利用其它的定时器，（2）通过忙循环法完成延时，或者（3）通过调用RMP\_Thd\_Delay完成延时。方法（3）是最节约处理器性能的方法，但并不总是可行，因为固件库可能在锁调度器或关中断期间调用延时；如果此时又没有其它定时器可用，则只能选择RMP\_Thd\_Loop提供的忙循环法完成延时。

定时器中断法的延时函数通过初始化定时器中断，然后在中断内递增某全局变量，且在延时时不断查询该全局变量来确定延时是否完成。这种延时函数的延时精度高、延时范围宽，而且还能够提供一个时间戳变量以供查询，但也可能干涉RMP的节拍定时器。不过，这种延时函数必然只能在中断开启时使用，因此可以用RMP\_Thd\_Delay进行简单替换。如果该延时函数产生的全局变量也被读取，则可以使用RMP\_Timestamp替代它。

\subsection{代码生成器的兼容性}\label{subsec:code-generator-compatibility}

为了使微控制器更加易用，部分厂商在库函数之上还额外提供了代码生成器。这些代码生成器的产出有可能与起冲突。常见的冲突有如下三种：

\begin{enumerate}
\def\labelenumi{\arabic{enumi}.}
\item
  生成含有函数的代码。为了兼容各厂商提供的启动代码对函数的调用，在内核中声明了函数。当代码生成器生成的框架与一起编译时，函数就会被定义两次，并在链接时引发错误。鉴于这些函数包含的大多是简单的硬件初始化代码，可以将该函数直接改名或宏定义为，然后将配置头文件中的定义成它。为了简化系统启动的心智模型\footnote{在函数进入后，部分存储器控制器、段寄存器等可能未初始化，且此时的执行并不处于线程环境中，进行某些内核操作可能会出错而且原因很难排查，而在一切都初始化好后再行操作则保证不会出现此类问题；此外，我们希望保证系统的入口总是函数，这样就能直接用它的栈来初始化中断响应栈，利于适配各裸机工具链；鉴于上述两点，我们不采用其它RTOS普遍使用的独立内核入口函数设计}，并不像那样提供独立的，因此其它代码需要避让函数。
\item
  生成初始化看门狗、开启某些无关外设中断等会意外干扰系统运行的代码。最好的解决方案是配置生成器使其不生成该种代码；如果这不可能，则必须手动修改生成后的代码，除去这些干扰项，或者在系统启动和运行过程中加入额外操作\footnote{如喂看门狗等}使它们不造成干扰。
\end{enumerate}

\section{RMP对C语言标准和MISRA标准的背离}\label{sec:rmp-deviations-from-the-c-standard-and-misra}

\subsection{C未定义行为：类型双关指针别名可以存在并不遵循严格别名机制}\label{subsec:c-undefined-behaviour-type-punning-pointer-aliasing-may-exist-and-does-not-follow-strict-aliasing}

类型双关（Type-Punning）指针别名（Pointer-Aliasing）是一种广泛存在于C语言\footnote{指C90，下同}实际应用中的用法，它允许两个不同类型的指针引用同一块内存。这在网络协议栈中是非常常见的，因为我们往往需要组装或者拆解网络封包，变换其字节序，等等。在操作系统中这也很常见，尤其是将无符号整数指针与结构体指针来回转换以方便对对象的操作。不幸的是，C语言的严格别名机制（Strict
Aliasing
Rule）并不允许两个不同类型的指针引用同一块内存，除非是兼容类型或者从其他类型转换到char\hyperref[sec:chapter-bibliography-11]{{[}1{]}}，而且也没有通用的办法绕开这一点。通常而言有三种方法可以避免类型双关，它们分别是：

\begin{enumerate}
\def\labelenumi{\arabic{enumi}.}
\item
  使用union来进行不同类别之间的转换。不幸的是这在C中是被定义行为而在C++中不是，而RMP的代码由于可能和使用C++写成的某些库直接链接\footnote{如果不希望使用extern "C"的话}而需要作为C++被编译，因此仍然有可能是未定义行为。此外，此种方法对长度在编译时不确定的数组对象是束手无策的，除非牺牲潜在的性能来对它们进行逐一转换。最后，这样使用union不符合union的设计精神，属于对语言规则的滥用。
\item
  使用char*来进行操作，因为char*总是可以和任何指针别名。不幸的是现存编译器对这一点的优化非常糟糕；某些相当优秀的编译器甚至也会把这些操作编译成逐字节访问的指令，对于低质量嵌入式编译器则更是如此。因此，这是一种不好的解决方案。
\item
  使用memcpy函数在两个不兼容的指针间拷贝内容，操作完后再拷贝回去，寄希望于编译器能够探测这种行为并且优化掉重复的memcpy。这巧妙地利用了memcpy的参数为char*的特点。它带来的直接问题是，某些低质量嵌入式编译器可能不够智能，有时不会把memcpy函数优化掉，而导致生成的代码中实际调用了两次memcpy，从而导致巨大的性能损失。而且一旦memcpy没有被优化掉，这可能会导致在函数内部声明的局部拷贝缓冲区被分配在栈上，大大增加栈消耗而导致堆栈溢出。最后，此类对memcpy的使用是不符合memcpy的设计意图的，属于对语言漏洞的滥用；C标准对此类做法不持鼓励态度。

  有鉴于以上三种方法的问题，RMP并不使用这些方法，而是直接使用类型双关指针别名，而在编译时关掉相关的优化\footnote{如GCC的-fno-strict-aliasing选项}。
\end{enumerate}

\subsection{C未定义行为：一个字节的长度总是八位，一个char类型的长度总是一个字节}\label{subsec:c-undefined-behaviour-a-byte-is-always-eight-bits-and-a-char-is-always-one-byte}

C语言并未定义char的长度，它仅仅表示char必须等于系统中基本字符集的单个字符的宽度\hyperref[sec:chapter-bibliography-11]{{[}1{]}}，理论上讲一个char类型可以是16或32个二进制位\footnote{某些早期数字信号处理器，如C28x的编译器会这样做；但现代数字信号处理器很少再这样做}，或者甚至更多。现时的所有主流编译器和架构上char的长度都总是8个二进制位，因此我们在代码中也作此假设以避免不必要的``char类型长度''配置宏。

事实上，RMP不会使用char型的高于8位的二进制位，也不会按字节操作或以字节为基础单位衡量更长的数据类型\footnote{使用的是按位操作和按char衡量，这在所有处理器上都是通行的}，因此在char型长度超过8位的前提下仍能正常运作\footnote{见C28x的移植}。不过，使用这类移植时必须小心：（1）rmp\_u8\_t不再是字面意义的8位，而是可以表达8位的最小整数类型\footnote{类似于stdint.h头文件中的uint\_least8\_t}，（2）本手册中所有的地址都是按字寻址，而且（3）本手册中所有的数据结构大小和内存长度都是指字数而非字节数。

\subsection{C未定义行为：指针的实际表示总是机器可直接识别的地址}\label{subsec:c-undefined-behaviour-a-pointer-actual-representation-is-always-a-machine-recognisable-address}

\protect\phantomsection\label{subsec:c-undefined-behaviour-a-pointer-actual-representation-is-always-a-machine-recognisable-address-2}C语言并未定义指针在内存中应该如何表示。理论上讲，指针变量的实际表示可以是任何东西，只要在解引用时它返回的值是按照C语言语义指向的变量的值就可以了\hyperref[sec:chapter-bibliography-11]{{[}1{]}}。这使得在严格的C标准下将指针对齐到某个内存粒度是不可能的，因为无法对其使用按位与来掩蔽掉它的某些位。因此，我们假设指针的实际表示总是机器可直接识别的二进制地址，这一点在所有RMP支持或计划支持的架构上都是成立的。

\subsection{C未定义行为：强制转换指针到不小于机器字长的无符号整数值总是不改变它的实际表示}\label{subsec:c-undefined-behaviour-converting-a-pointer-to-an-unsigned-integer-no-smaller-than-a-machine-word-always-preserves-its-actual-representation}

C语言并未规定指针在转换到无符号整数后，其实际表示应如何变化。它只规定了，如果先把一个指针转换成一个空间足够大的无符号整数后再转换回去，来回转换后得到的指针在和原指针做比较操作时应当视为相等\footnote{这甚至不禁止转换前和来回转换后得到的表示有差异，只要这两个值在做指针比较时总是按照相等处理即可}\hyperref[sec:chapter-bibliography-11]{{[}1{]}}。如果指针在转换到无符号整数后的表示会变化，那么我们会遇到与\ref{subsec:c-undefined-behaviour-a-pointer-actual-representation-is-always-a-machine-recognisable-address}
所述类似的问题。因此，我们假设强制转换指针到一个机器字长的无符号整数值总是不改变它的实际表示，这一点在所有RMP支持或计划支持的架构上都是成立的。

事实上，RMP并不禁止在转换到更长的无符号整数时在该整数的长于原指针的高位添加任何内容\footnote{如RL78的CC-RL编译器会在内存指针转换得到的无符号整数的高位插入0x000F，代表这是指向内存空间}，只要这些高位内容在转换回原指针类型后被丢弃；这是因为，RMP只会对该无符号整数的最低几位进行掩蔽操作，然后再将该无符号整数转换成指向同一类别存储器空间\footnote{部分架构的数据和代码存储器是分离的；数据和代码存储器也可能各有多种子类别}的（1）原指针类型或（2）不比原指针类型的内存对齐条件更苛刻的指针类型使用，不会涉及到潜在的高位内容。不过，在使用此种移植时必须注意，由指针转换成的无符号整数会具备意外的高字节垃圾。

\subsection{C未定义行为：一切``严丝合缝''的结构体会自动紧凑排列}\label{subsec:c-undefined-behaviour-all-exactly-fitting-structures-are-automatically-packed}

C语言并未规定结构体中各变量的内存地址对齐方式。理论上讲，编译器可以任意决定结构体中的变量的对齐方式，只要它保证结构体的访问语义就可以了\hyperref[sec:chapter-bibliography-11]{{[}1{]}}。这甚至不禁止将一个仅含有一种类型的结构体对齐到更大的任意粒度\footnote{比如将仅含char类型的结构体中的每个变量对齐到1GiB；这种荒唐行为是C语义容许的}。因此，即便结构体在类型方面是``纯洁''的，我们也无法保证它内部的变量是紧凑排列的。为解决C语言规范在这方面的缺失，几乎所有编译器都提供了额外机制来使结构体中各变量紧凑排列\footnote{其中最出名的是GCC的\_\_attribute\_\_((packed))}，但它们提供的手段堪比巴别塔：有的编译器使用\#pragma，有的编译器使用\_\_attribute\_\_，而且不同编译器的机制在结构体声明中的位置是不同的。使用\#pragma的编译器尤其棘手，因为宏无法展开成\#pragma；该\#pragma必须被直接插入代码文件中。

为了避免上述问题，RMP假设编译器总是自动紧凑排列一切``严丝合缝''的结构体，其中``严丝合缝''的定义如下：

\begin{enumerate}
\def\labelenumi{\arabic{enumi}.}
\item
  结构体中的``最终基础类型''的大小都是2的次方；其中，``最终基础类型''是指递归组成该结构体类型的编译器原生数据类型；
\item
  若结构体被紧凑排列，则其中任何``最终基础类型''的变量的偏移地址都自然对齐到该``最终基础类型''的大小；
\item
  结构体的基地址对齐要求不高于其中最长的``最终基础类型''的对齐要求。

  几乎所有的编译器都满足上述要求，因为对齐变量的唯一目的是防止发生可能影响性能或造成异常的非对齐访问。有鉴于此，RMP不保证支持不满足这个要求的编译器。
\end{enumerate}

\subsection{C未定义行为：外部标识符的长度扩展为31字符}\label{subsec:c-undefined-behaviour-external-identifiers-are-extended-to-31-characters}

C90规定外部（extern）变量名最少只有前6个字符是有效的。这对任何操作系统来说都太少了。因此，RMP把这个限制放宽到了与内部变量名限制相同的31个字符。所有的现代编译器均支持这一点。

\subsection{MISRA-C背离情况说明}\label{subsec:misra-c-deviations}

RMP对于MISRA-C:2012\hyperref[sec:chapter-bibliography-11]{{[}2{]}}中的所有三类规则都进行了严格检查，并且所有的背离都被批准。

RMP对于MISRA-C:2012\hyperref[sec:chapter-bibliography-11]{{[}2{]}}中所述的强制（Mandatory）类别没有任何背离发生。

RMP对于MISRA-C:2012\hyperref[sec:chapter-bibliography-11]{{[}2{]}}中所述的要求（Required）类别的背离如下：

\captionof{table}{RMP对MISRA-C:2012中所述的要求（Required）部分的背离}\label{tab:rmp-deviations-from-the-required-part-of-misra-c-2012}

{\def\LTcaptype{none} % do not increment counter
\begin{longtable}[]{@{}
  >{\raggedright\arraybackslash}p{(\linewidth - 2\tabcolsep) * \real{0.0841}}
  >{\raggedright\arraybackslash}p{(\linewidth - 2\tabcolsep) * \real{0.9159}}@{}}
\toprule\noalign{}
\endhead
\bottomrule\noalign{}
\endlastfoot
条目 & 解释 \\
1.3 &
某些如前所述的未定义行为对于实现操作系统这样的底层软件而言是必需的。 \\
5.1 &
RMP扩展了外部标识符的有效长度到31位。这对于操作系统的软件工程而言是必需的。 \\
8.4 &
RMP使用特殊的头文件编写方法规避了这一问题，并使得变量的声明和定义可以在一起进行。这有利于提高软件质量。 \\
11.3

11.6 &
类型双关指针别名和指针到足够长度的整数的转换对于操作系统这样的软件而言是必需的。RMP事实上仅在（1）结构体和子结构体处、（2）HAL接口处和（3）内存池处使用此类双关指针别名。 \\
11.9 &
为防止与其它第三方源码包冲突，RMP不使用NULL而使用RMP\_NULL来表示空指针。 \\
14.3 &
RMP在检查配置是否正确的流程中大量使用了编译时就能确定结果的布尔表达式。这些布尔表达式都被直接送入RMP\_ASSERT宏，确保RMP被正确配置。这提高了软件的质量。 \\
17.7 &
由于RMP的硬件抽象层要对所有的受支持硬件定义统一的接口，因此不免会有一些硬件抽象层或/和调试打印函数的返回值对于某些架构而言是多余的。 \\
21.1 &
RMP声明了\_\_HDR\_DEF\_\_、\_\_HDR\_STRUCT\_\_和\_\_HDR\_PUBLIC\_\_三个保留宏，它们仅做自用。因此，这并不违反\#undef系统中其他保留宏的规定。另见20.5的说明。 \\
\end{longtable}
}

RMP对于MISRA-C:2012\hyperref[sec:chapter-bibliography-11]{{[}2{]}}中所述建议（Advisory）类别的背离如下：

\captionof{table}{RMP对MISRA-C:2012中所述的建议（Advisory）部分的背离}\label{tab:rmp-deviations-from-the-advisory-part-of-misra-c-2012}

{\def\LTcaptype{none} % do not increment counter
\begin{longtable}[]{@{}
  >{\raggedright\arraybackslash}p{(\linewidth - 2\tabcolsep) * \real{0.0841}}
  >{\raggedright\arraybackslash}p{(\linewidth - 2\tabcolsep) * \real{0.9159}}@{}}
\toprule\noalign{}
\endhead
\bottomrule\noalign{}
\endlastfoot
条目 & 解释 \\
2.3

2.5 &
在硬件抽象层，尤其是固件库中，某些宏或数据类型是打包一起提供的。它们可能在当前配置中没有使用，但必须保留以适应未来的软件配置调整，有利于维护工程的一致性和提高代码质量。这些宏都有很好的注释和文档可供参考。 \\
2.7

8.7 &
由于RMP的硬件抽象层要对所有的受支持硬件定义统一的接口，因此不免会有一些硬件抽象层函数的参数对于某些架构而言是多余的，或者部分声明了extern的函数未被实际调用，等等。 \\
8.9 &
RMP有部分全局变量仅被一个函数引用，但并不放在该函数内部，这是为了方便用户进行代码调试。 \\
11.4

11.5 &
转换一个无符号数或者void*到常规指针对于操作系统这样的软件而言是必需的。具体解释见11.3的相关说明。 \\
15.5 &
RMP的一个函数可能会有多个出口。在这些代码处都进行了良好注释，并且其他可选实现会违反其他的要求或建议，并且破坏工程的整体抽象。 \\
20.5 &
为了使每个函数和变量的声明都只进行一次（包括其extern形式），系统性地防止出现隐式声明错误来提高软件质量，RMP内部使用了特殊的分段包含式头文件。该格式要求使用\#undef。\#undef仅用于该内部格式，不会暴露到内核外部，且该格式在整个系统中都被良好遵循，并有对应文档进行详细说明。 \\
\end{longtable}
}

\section{本章参考文献}\label{sec:chapter-bibliography-10}

\protect\phantomsection\label{sec:chapter-bibliography-11}{[}1{]}
ISO, International Standards Organisation, "Programming Languages---C,
International standard, ISO/IEC 9899: 1990 (E)", 1990.

{[}2{]} M. I. S. R. Association, MISRA C 2012: Guidelines for the Use of
the C Language in Critical Systems: March 2013: Motor Industry Research
Association, 2013.
